\subsection{Rastreabilidade entre Requisitos de Sistema (SRS) e FRETish}
A especificação dos requisitos de engenharia do sistema de telemetria (\textit{System Requirements Specification} -- SRS) foi integralmente formalizada na gramática FRETish. A Tabela~\ref{tab:srs_fret} apresenta a matriz de rastreabilidade exaustiva entre os requisitos de sistema (\texttt{REQ-SYS-01} a \texttt{REQ-SYS-30}), os 9 componentes arquiteturais de software e emulação HIL, os identificadores formais no FRET e o resultado de realizabilidade obtido no provador formal Kind~2 / Z3.

\begin{table}[htbp]
\centering
\scriptsize
\caption{Matriz de Rastreabilidade Completa: Requisitos de Sistema (SRS) vs. Requisitos Formais FRET (Pilar~1).}
\label{tab:srs_fret}
\setlength{\tabcolsep}{3.5pt}
\renewcommand{\arraystretch}{1.18}
\begin{tabularx}{\textwidth}{@{} l >{\raggedright\arraybackslash}X l >{\raggedright\arraybackslash}p{3.2cm} c c @{}}
\toprule
\textbf{ID SRS} & \textbf{Descrição do Requisito de Sistema (SRS)} & \textbf{Componente} & \textbf{Requisitos FRETish} & \textbf{Critério} & \textbf{Status Kind 2} \\
\midrule
\texttt{REQ-SYS-01, 02} & Emissão periódica DBC e dinâmica do motor HIL & \texttt{uno\_ecu\_emulator} & \texttt{REQ\_EMU\_001..005} & [AC-01, 03] & \textbf{Realizable} \\
\texttt{REQ-SYS-03, 04} & Ingestão assíncrona ISR e taxa de perda $\le 1\%$ & \texttt{esp32\_twai} & \texttt{REQ\_CAN\_001..006} & [AC-01] & \textbf{Realizable} \\
\texttt{REQ-SYS-05, 06} & Polling cíclico a 10~Hz e escalonamento \textit{Round-Robin} & \texttt{esp32\_obd} & \texttt{REQ\_OBD\_001, 002, 004} & [AC-02] & \textbf{Realizable} \\
\texttt{REQ-SYS-07..09} & Emissão de respostas OBD-II com jitter $< 3$~ms & \texttt{uno\_ecu\_emulator} & \texttt{REQ\_EMU\_006..008} & [AC-02, 03] & \textbf{Realizable} \\
\texttt{REQ-SYS-10}     & Detecção e tratamento de timeout de resposta (50~ms) & \texttt{esp32\_obd} & \texttt{REQ\_OBD\_003} & -- & \textbf{Realizable} \\
\texttt{REQ-SYS-11..13} & Serialização CSV sem alocação dinâmica e buffer SRAM & \texttt{esp32\_logger} & \texttt{REQ\_LOG\_001..003} & [AC-04] & \textbf{Realizable} \\
\texttt{REQ-SYS-14, 15} & \textit{Flush} híbrido por limiar (3584B) / tempo (2s) e DMA & \texttt{esp32\_logger} & \texttt{REQ\_LOG\_004..006} & [AC-04] & \textbf{Realizable} \\
\texttt{REQ-SYS-16, 17} & Montagem FAT32 no SPI2 e rotação atômica de arquivos & \texttt{esp32\_sd} & \texttt{REQ\_SD\_001..003} & -- & \textbf{Realizable} \\
\texttt{REQ-SYS-18}     & Encerramento seguro de sessão e volume $\ge 72.000$ amostras & \texttt{esp32\_sd} & \texttt{REQ\_SD\_004, 005} & [AC-08] & \textbf{Realizable} \\
\texttt{REQ-SYS-19..22} & Conexão Wi-Fi, publicação MQTT e despacho no ciclo 50~ms & \texttt{esp32\_telemetry} & \texttt{REQ\_COM\_001..005} & [AC-04] & \textbf{Realizable} \\
\texttt{REQ-SYS-23, 24} & \textit{Fallback} determinístico SD ($\le 5$~ms) e dreno de backlog & \texttt{esp32\_fsm, logger} & \texttt{REQ\_FSM\_001, 002}, \texttt{007} & [AC-06] & \textbf{Realizable} \\
\texttt{REQ-SYS-25, 26} & Validação e roteamento de comandos remotos MQTT e replay & \texttt{esp32\_cmd} & \texttt{REQ\_CMD\_001..003} & -- & \textbf{Realizable} \\
\texttt{REQ-SYS-27}     & Auto-recuperação de Bus-Off com handshake ISO~11898 & \texttt{esp32\_recovery} & \texttt{REQ\_REC\_001..005} & [AC-07] & \textbf{Realizable} \\
\texttt{REQ-SYS-28, 30} & Detecção de \textit{stall} do logger (30s) e watchdog hardware & \texttt{esp32\_recovery} & \texttt{REQ\_REC\_006, 007} & -- & \textbf{Realizable} \\
\texttt{REQ-SYS-29}     & Teto estático de SRAM $< 200$~KB e batimento de integridade & \texttt{esp32\_logger} & \texttt{REQ\_LOG\_008, 009} & [AC-05] & \textbf{Realizable} \\
\bottomrule
\end{tabularx}
\end{table}


\subsection{Rastreabilidade com os Critérios de Aceitação da Dissertação (AC-01 a AC-08)}
A engenharia de requisitos do coletor vincula rigorosamente os critérios de aceitação empíricos da dissertação aos contratos matemáticos do sistema, conforme sumarizado na Tabela~\ref{tab:rastreabilidade_ac}.

\begin{table}[htbp]
\centering
\footnotesize
\caption{Matriz de Rastreabilidade: Critérios de Aceitação da Dissertação vs. Requisitos FRET.}
\label{tab:rastreabilidade_ac}
\setlength{\tabcolsep}{4pt}
\renewcommand{\arraystretch}{1.25}
\begin{tabular}{@{} c >{\raggedright\arraybackslash}p{6.6cm} >{\raggedright\arraybackslash}p{4.5cm} >{\raggedright\arraybackslash}p{2.6cm} @{}}
\toprule
\textbf{Critério} & \textbf{Descrição do Alvo de Engenharia} & \textbf{Requisitos FRETish} & \textbf{Componente} \\
\midrule
\textbf{[AC-01]} & Taxa de recepção e perda de tráfego CAN $\le 1{,}0\%$ a 500~kbps & \texttt{REQ\_CAN\_001, 002, 005} & \texttt{esp32\_twai} \\
\textbf{[AC-02]} & \textit{Polling} ativo cíclico e latência média OBD-II $< 10{,}0$~ms & \texttt{REQ\_OBD\_001, 002}, \newline \texttt{REQ\_EMU\_006} & \texttt{esp32\_obd}, \newline \texttt{uno\_ecu} \\
\textbf{[AC-03]} & Estabilidade temporal e Jitter de resposta OBD-II $< 3{,}0$~ms & \texttt{REQ\_EMU\_005, 008} & \texttt{uno\_ecu\_emulator} \\
\textbf{[AC-04]} & Integridade da serialização CSV e \textit{Throughput} $\ge 200$~pkt/s & \texttt{REQ\_LOG\_001, 003}, \newline \texttt{REQ\_COM\_002} & \texttt{esp32\_logger}, \newline \texttt{telemetry} \\
\textbf{[AC-05]} & Consumo estático de memória SRAM $< 200{,}0$~KB & \texttt{REQ\_LOG\_008, 009} & \texttt{esp32\_logger} \\
\textbf{[AC-06]} & \textit{Fallback} offline no MicroSD sob falha de rede $\le 5{,}0$~ms & \texttt{REQ\_FSM\_001, 002}, \newline \texttt{REQ\_COM\_003} & \texttt{esp32\_fsm}, \newline \texttt{telemetry} \\
\textbf{[AC-07]} & Recuperação automática de Bus-Off com espera de 128~ms (ISO~11898) & \texttt{REQ\_REC\_001} a \texttt{005} & \texttt{esp32\_recovery} \\
\textbf{[AC-08]} & Gravação contínua no SD $\ge 72.000$~amostras em ensaio de 1 hora & \texttt{REQ\_SD\_001, 004, 005} & \texttt{esp32\_sd} \\
\bottomrule
\end{tabular}
\end{table}


\subsection{Metodologia Baseada em Sheridan, Becker et al. (RefSQ 2025)}
A formalização em FRETish seguiu as diretrizes e lições aprendidas documentadas pelo grupo de pesquisa do Prof.~Dr.~Leandro Buss Becker (UFSC) na conferência RefSQ~2025 \cite{sheridan2025}. A experiência prática de desenvolvimento revelou que a formalização de requisitos não constitui mera tradução textual, mas um processo evolutivo de refinamento de engenharia:
\begin{enumerate}
    \item \textbf{Tentativa Ingênua Inicial (Monolítica):} Em um primeiro momento, concentraram-se mais de 40 requisitos em um único componente monolítico (\texttt{esp32s3\_collector}). Prazos físicos longos foram declarados explicitamente na sintaxe temporal (ex.: temporização de $128\text{ ms}$ no Bus-Off como \texttt{shall after} \allowbreak \texttt{128 MILLISECOND} e \textit{timeout} de rede de $10\text{ s}$). O compilador Lustre realizou o desenrolamento (\textit{unrolling}) discreto de cada milissegundo em variáveis de estado interno, gerando um modelo com mais de 20.000 linhas ($\approx 1{,}5\text{ MB}$). O provador formal SMT Kind~2 esgotou a memória RAM do sistema operacional (\textit{Out-of-Memory}), retornando \texttt{UNKNOWN / Timeout} após 15 minutos;
    \item \textbf{Diagnóstico e Decomposição em Componentes de Software (SW-C):} Apoiando-se no precedente de Becker et al. \cite{sheridan2025} e nas normas NASA-STD-8739.8, a arquitetura foi decomposta em 8 componentes de software especializados (\texttt{esp32\_twai}, \texttt{esp32\_obd}, \texttt{esp32\_logger}, \texttt{esp32\_sd}, \allowbreak \texttt{esp32\_telemetry}, \texttt{esp32\_fsm}, \texttt{esp32\_cmd}, \texttt{esp32\_recovery}) e 1 componente emulador de bancada (\texttt{uno\_ecu\_emulator});
    \item \textbf{Adoção do NASA Timer Handshake Pattern:} Temporizações físicas longas de processos externos foram desacopladas dos prazos estritos de CPU. O software dispara o temporizador com o predicado (\texttt{immediately satisfy} \allowbreak \texttt{recovery\_timer\_128ms\_start}) e processa em prazo delimitado ($\le 1\text{ ms}$) o recebimento do término (\texttt{upon} \allowbreak \texttt{recovery\_timer\_128ms\_expired});
    \item \textbf{Garantia Matemática de Realizabilidade:} Com a reestruturação, o modelo Lustre compilado reduziu-se a 5~KB. A verificação formal monolítica e composicional no Kind~2 / Z3 convergiu em menos de \textbf{0{,}41~segundos}, alcançando \textbf{100\% de aprovação (\texttt{Realizable: True})} para todos os 48 requisitos formais, sem inconsistências lógicas ou \textit{deadlocks}.
\end{enumerate}

\begin{figure}[htbp]
    \centering
    \begin{subfigure}[b]{0.48\textwidth}
        \centering
        \includegraphicssafe{width=\linewidth}{fret_dashboard.png}{Dashboard do FRET comprovando 48 requisitos e 100\% formalizados.}
        \caption{Painel geral do NASA FRET (48 requisitos em 9 componentes).}
        \label{fig:fret_dashboard}
    \end{subfigure}
    \hfill
    \begin{subfigure}[b]{0.48\textwidth}
        \centering
        \includegraphicssafe{width=\linewidth}{fret_realizability.png}{Tela do FRET executando Realizability Checking via Kind 2.}
        \caption{Verificação de Realizabilidade via Kind~2 / SMT.}
        \label{fig:fret_realizability}
    \end{subfigure}
    \caption{Evidências de formalização e verificação de realizabilidade no NASA FRET.}
    \label{fig:fret_evidencias}
\end{figure}

\subsection{Estrutura e Amostra de Sentenças FRETish}
Cada sentença formalizada em FRETish decompõe-se rigorosamente em cinco cláusulas semânticas: \textbf{Scope} (fronteira operacional), \textbf{Condition} (gatilho de ativação), \textbf{Component} (agente responsável), \textbf{Timing} (restrição temporal estrita) e \textbf{Response} (predicado lógico pós-condição). O catálogo completo reúne 48 requisitos; a Listagem~\ref{lst:fret_amostra} ilustra as sentenças representativas dos componentes críticos.

\begin{lstlisting}[language=FRETish, caption={Amostra de sentenças representativas codificadas em FRETish.}, label={lst:fret_amostra}]
-- REQ_CAN_002: Ingestao deterministica na ISR de recepcao TWAI [AC-01]
in active_session upon can_frame_arrived the esp32_twai shall within 2 MILLISECOND satisfy frame_timestamp_captured

-- REQ_OBD_001: Disparo de polling ciclico SAE J1979 a 10 Hz [AC-02]
in active_session upon obd_poll_timer_100ms_tick the esp32_obd shall within 1 MILLISECOND satisfy obd_request_transmitted & current_pid_rotated

-- REQ_LOG_001: Formatacao CSV em buffer estatico de SRAM [AC-04, AC-05]
in active_session upon telemetry_sample_ready the esp32_logger shall within 1 MILLISECOND satisfy csv_line_formatted & no_heap_allocation

-- REQ_FSM_001: Comutacao automatica para MicroSD em caso de perda de Wi-Fi [AC-06]
in active_session upon wifi_disconnected the esp32_fsm shall within 5 MILLISECOND satisfy fallback_to_sd_active

-- REQ_REC_002: Handshake do temporizador ISO 11898 de 128 ms para Bus-Off [AC-07]
in bus_off_state upon bus_off_detected the esp32_recovery shall immediately satisfy recovery_timer_128ms_start
\end{lstlisting}

\subsection{Prontidão para Runtime Verification (RV) com R2U2 e Lógica MLTL}
Conforme destacado na metodologia de Becker et al. \cite{sheridan2025}, a formalização em FRETish viabilizou a exportação direta de contratos temporais para Verificação em Tempo de Execução (\textit{Runtime Verification} -- RV) através do motor R2U2 (\textit{Realizable, Responsive, Unobtrusive Unit}). Todos os 48 requisitos geraram fórmulas compiladas em Lógica Temporal de Tempo de Missão (\textit{Mission-time Linear Temporal Logic} -- MLTL) com janelas delimitadas de clock, permitindo que oráculos dinâmicos monitorem o barramento CAN em bancada sem sobrecarga computacional.
